%% ****** Start of file main_pra.tex ******
%%
%% REVTeX 4.2 manuscript for Physical Review A (Regular Article)
%%
%% Title: Canonical Conjugate Dynamics of the Fermionic Correlation Matrix
%%        and the Mode Decomposition Necessity of Complex Numbers in Quantum Mechanics
%%
%% Compiled with:
%%   pdflatex main_pra
%%   bibtex  main_pra
%%   pdflatex main_pra
%%   pdflatex main_pra
%%

\documentclass[aps,pra,preprint,groupedaddress,amsmath,amssymb,floatfix]{revtex4-2}

\usepackage{graphicx}
\usepackage{bm}
\usepackage{hyperref}

\begin{document}

\title{Canonical Conjugate Dynamics of the Fermionic Correlation Matrix
       and the Mode Decomposition Necessity of Complex Numbers in Quantum Mechanics}

\author{[Author names]}
\affiliation{[Institution]}
\email{[email]}

\date{\today}

\begin{abstract}
Quantum mechanics rests on two postulates---the Schr\"odinger equation (unitary, complex-valued, reversible) and the Born rule (probabilistic, real-valued, irreversible)---whose structural tension has remained unresolved for a century. This ``Born--Schr\"odinger gap'' is widely regarded as a conceptual puzzle. We show that it is instead a physical necessity arising from the canonical conjugate structure of the fermionic correlation matrix $C_{mn} = \langle c^\dagger_n c_m \rangle$. Starting from the Heisenberg equation of motion for $C$, we prove that its real part $C^R$ and imaginary part $C^I$ evolve as exact canonical conjugates (Identity~1), with dynamics that preserve the symplectic $2$-form $\Omega = \sum_{k<l} dC^R_{kl} \wedge dC^I_{kl}$. The real evolution matrix for each conjugate pair is $M = \bigl(\begin{smallmatrix}0 & \omega \\ -\omega & 0\end{smallmatrix}\bigr)$, whose eigenvalues $\pm i\omega$ are purely imaginary---this matrix is not diagonalizable over $\mathbb{R}$. We prove that this algebraic incompleteness forces the minimal sufficient description to be over $\mathbb{C}$ (Theorem~1): the complex structure of the correlation matrix dynamics is dynamically irreducible and cannot be reduced to a real-valued description. Building on this foundation, we derive an exact decomposition of the ergotropy rate (Theorem~2):
\[
\dot{\mathcal{E}} = -\sum_{m\neq n} \omega_{mn} (\Delta H_\rho)_{nm} C^I_{mn},
\]
showing that the imaginary part of the correlation matrix---precisely the information discarded by the Born rule projection $|\cdot|^2$---is the unique driver of extractable work in quantum systems. We propose an operational weak measurement protocol for $C^I$ that provides an independent experimental path to this quantity, circumventing the Born rule projection on the system. Two testable predictions follow: (A)~a state-dependent uncertainty relation $\Delta C^R \cdot \Delta C^I \ge \frac{1}{4}|\langle n_m\rangle - \langle n_n\rangle|$, whose lower bound is set by a directly measurable density gradient, offering a sharp experimental signature in quantum gas microscopes; and (B)~a closed-form power formula for a three-level single-spin-device quantum engine, verified numerically over $405$ parameter combinations. These results reframe the Born--Schr\"odinger gap from a conceptual puzzle into a productive physical mechanism, identifying complex-valued coherence as the operational fuel of quantum thermodynamic processes.
\end{abstract}

\maketitle

%% =================================================================
%% Section I: Introduction
%% =================================================================
\section{Introduction}
\label{sec:introduction}

Quantum mechanics was assembled from two pieces that have never quite fit together. The Schr\"odinger equation prescribes unitary, deterministic evolution in a complex vector space; the Born rule delivers real-valued probabilities upon measurement. The former is reversible and phase-coherent; the latter is irreversible and discards phase. This ``Born--Schr\"odinger gap'' has been noted since the founding of the theory, but its physical status---accidental feature or necessary structural consequence---has remained unclear.

A necessary logical clarification is in order before proceeding. Our logical structure is:

\begin{enumerate}
\item[(i)] The fermionic correlation matrix $C_{mn} = \langle c^\dagger_n c_m \rangle$ is an experimentally measurable quantity whose complex nature follows from the canonical anti-commutation relations~--- this is an observational fact, not a theoretical assumption.
\item[(ii)] Given that $C$ is complex-valued, we decompose $C = C^R + i C^I$.
\item[(iii)] The real and imaginary parts evolve under canonical conjugate dynamics (Identity~1).
\item[(iv)] We ask: can the dynamics of $(C^R,C^I)$ be consistently described using only real numbers?
\item[(v)] We prove the answer is no~--- the evolution matrix is not diagonalizable over $\mathbb{R}$ (Theorem~1).
\item[(vi)] Therefore, the complex structure at the level of the correlation matrix is dynamically irreducible~--- it cannot be reduced to a real-valued description.
\end{enumerate}

The reader may object that the $2\times 2$ matrix $M = \bigl(\begin{smallmatrix}0 & \omega \\ -\omega & 0\end{smallmatrix}\bigr)$ also describes the classical harmonic oscillator, which is perfectly well represented over $\mathbb{R}$. This objection misunderstands the role of spectral decomposition in quantum versus classical physics. In classical mechanics, the generator $M$ is never diagonalized---the second-order ODE $m\ddot{x} + kx = 0$ is solved directly in the time domain, yielding real sinusoids. Classical physics has no need for eigenstates, quantum numbers, or spectral decomposition, because the system's identity is given by its trajectory, not by its Hamiltonian eigenbasis. In quantum mechanics, by contrast, the Hamiltonian eigenbasis defines the physical language of the theory: energy levels, occupation numbers, Fock space, and quantum numbers are all defined through spectral decomposition. This spectral demand forces the transition to $\mathbb{C}$, because $M$ over $\mathbb{R}$ has no eigenvectors---the dynamics is permanently coupled and cannot be decomposed into independent modes without complexifying. The classical harmonic oscillator does not need $\mathbb{C}$ precisely because it does not need spectral decomposition. Thus the apparent counterexample is, in fact, a precise articulation of the quantum-classical divide.

We do not claim to derive complex numbers from real numbers alone; we show that the experimentally-observed complex correlation matrix has a dynamical structure that prohibits consistent real-valued reduction.

\textbf{Existing approaches.} The question of why quantum mechanics is complex has been approached from several directions. Kibble~\cite{kibble1979} and Ashtekar and Schilling~\cite{ashtekar1999} showed that the geometric structure of complex projective space unifies the Schr\"odinger equation and the Born rule, but their framework assumes complex amplitudes from the outset and does not address their dynamical origin. Renou~\textit{et al.}~\cite{renou2021} demonstrated that real quantum mechanics makes operationally different statistical predictions from complex quantum mechanics in network Bell experiments~[\textit{Nature} \textbf{600}, 625 (2021)]; their result establishes that real and complex quantum theories are empirically distinguishable, but does not trace the distinguishability to a dynamical root. de Oliveira~\cite{deoliveira2023} showed that canonical variables can be complexified, but concluded that complex numbers are ``natural'' rather than ``necessary''~[\textit{Braz. J. Phys.} \textbf{55}, 13 (2025)].

Most closely related in spirit is the work of Goyal, Knuth, and Skilling~\cite{goyal2010}, who derived complex quantum amplitudes from operational symmetry principles (Feynman's probability rules combined with continuous reversible transformations). Their derivation establishes the necessity of complex numbers at the level of measurement postulates. Our work identifies a different and complementary mechanism: at the level of the correlation matrix dynamics, the canonical conjugate structure of $C^R$ and $C^I$ forces an escape from $\mathbb{R}$ to $\mathbb{C}$ for purely algebraic reasons~--- the real evolution matrix is not diagonalizable. While Goyal~\textit{et al.} answer ``why complex amplitudes'' from operational axioms, we answer ``why complex correlation dynamics'' from the structure of the Heisenberg equation. The two derivations are independent and mutually reinforcing.

\textbf{This paper.} In Sec.~\ref{sec:identity1} we establish Identity~1: the real and imaginary parts of the fermionic correlation matrix evolve as canonical conjugates, with dynamics that preserve a symplectic $2$-form. Section~\ref{sec:theorem1} proves Theorem~1: the dynamical necessity of complex numbers follows from the non-diagonalizability of the real evolution matrix over $\mathbb{R}$. Section~\ref{sec:theorem2} derives Theorem~2: the ergotropy rate decomposes into a form driven exclusively by $C^I$, with an operational weak measurement protocol providing independent experimental access to the imaginary sector. Predictions derived from these results, together with their proposed experimental tests, are summarized.

%% =================================================================
%% Section II: Identity 1
%% =================================================================
\section{Identity 1: Canonical Conjugate Dynamics}
\label{sec:identity1}

%% -----------------------------------------------------------------
\subsection{Setup and notation}
\label{sec:setup}

Consider a fermionic system described by creation and annihilation operators $c^\dagger_m$, $c_m$ satisfying the canonical anti-commutation relations $\{c_m, c^\dagger_n\} = \delta_{mn}$. The one-body correlation matrix (also called the single-particle density matrix) is defined as

\begin{equation}
C_{mn} = \langle c^\dagger_n c_m \rangle,
\label{eq:C_def}
\end{equation}

where the expectation value is taken in the state of interest. $C$ is Hermitian by construction: $C^\dagger = C$.

We decompose $C$ into its real and imaginary parts,

\begin{equation}
C = C^R + i C^I,
\qquad
C^R = \frac{C + C^\dagger}{2},\quad
C^I = \frac{C - C^\dagger}{2i}.
\label{eq:C_decomp}
\end{equation}

Both $C^R$ and $C^I$ are real symmetric matrices. For a system with $N$ sites, $C$ has $N(N-1)/2$ independent complex off-diagonal elements; equivalently, $(C^R, C^I)$ constitutes $N(N-1)$ real degrees of freedom.

The Hamiltonian $H$ is taken to be real symmetric, which is the generic case for fermionic lattice models without magnetic fields.

%% -----------------------------------------------------------------
\subsection{Derivation of conjugate dynamics}
\label{sec:derivation}

The Heisenberg equation of motion for $C$ is

\begin{equation}
\dot{C} = i[H, C].
\label{eq:heisenberg}
\end{equation}

Substituting $C = C^R + i C^I$ and using the reality of $H$, $C^R$, and $C^I$:

\begin{equation}
\dot{C}^R + i\dot{C}^I = i[H, C^R + iC^I] = i[H, C^R] - [H, C^I].
\label{eq:heisenberg_expand}
\end{equation}

Separating real and imaginary parts gives the exact pair

\begin{equation}
\boxed{\frac{d}{dt} C^R = -[H, C^I], \qquad
\frac{d}{dt} C^I = +[H, C^R]}.
\label{eq:conjugate_dynamics}
\end{equation}

These equations are an algebraic identity following from the definition of $C$ and the Heisenberg equation. No approximation has been made; they hold for any state and any real symmetric $H$.

%% -----------------------------------------------------------------
\subsection{Symplectic structure}
\label{sec:symplectic}

Equation~(\ref{eq:conjugate_dynamics}) reveals that $(C^R, C^I)$ form canonical conjugate pairs. To make this structure precise, define the symplectic $2$-form on the space of correlation matrix elements:

\begin{equation}
\Omega = \sum_{k<l} dC^R_{kl} \wedge dC^I_{kl}.
\label{eq:symplectic}
\end{equation}

The evolution generated by Eq.~(\ref{eq:conjugate_dynamics}) preserves $\Omega$. To see this, compute the Lie derivative of $\Omega$ along the flow:

\begin{equation}
\mathcal{L}_{\dot{C}} \Omega = \sum_{k<l} \bigl( d\dot{C}^R_{kl} \wedge dC^I_{kl} + dC^R_{kl} \wedge d\dot{C}^I_{kl} \bigr).
\label{eq:lie_deriv}
\end{equation}

Using Eq.~(\ref{eq:conjugate_dynamics}), expand $\dot C^R_{ij} = -[H, C^I]_{ij}$ and $\dot C^I_{ij} = [H, C^R]_{ij}$. Taking exterior derivatives:
\begin{align*}
d(\dot C^R_{ij}) &= -\sum_{k} \bigl(H_{ik}\,dC^I_{kj} - H_{kj}\,dC^I_{ik}\bigr),\\
d(\dot C^I_{ij}) &= \;\;\,\sum_{k} \bigl(H_{ik}\,dC^R_{kj} - H_{kj}\,dC^R_{ik}\bigr).
\end{align*}
Substituting into Eq.~\eqref{eq:lie_deriv} and discarding $dC^I\wedge dC^I=0$ terms,
\begin{align*}
\mathcal{L}_{\dot C}\Omega = \sum_{i<j}\sum_{k} \bigl[&H_{kj}\,dC^I_{ik}\wedge dC^I_{ij}
- H_{ik}\,dC^I_{kj}\wedge dC^I_{ij} \\
&+ H_{ik}\,dC^R_{ij}\wedge dC^R_{kj}
- H_{kj}\,dC^R_{ij}\wedge dC^R_{ik}\bigr].
\end{align*}
The $dC^I$ terms vanish because $H_{kj}$ is symmetric while $dC^I_{ik}\wedge dC^I_{ij}$ is antisymmetric under $(k\leftrightarrow j)$; the $dC^R$ terms vanish identically. Hence
\begin{equation}
\mathcal{L}_{\dot C}\Omega = 0 \quad \Longrightarrow \quad \frac{d\Omega}{dt} = 0.
\label{eq:omega_conserved}
\end{equation}

The conservation of $\Omega$ means that $(C^R, C^I)$ parameterize a symplectic manifold whose symplectic volume is preserved by the Heisenberg evolution. Each pair $(C^R_{kl}, C^I_{kl})$ behaves as a canonical coordinate and momentum, and the generator $H$ acts as a Hamiltonian function on this phase space via the commutator.

%% -----------------------------------------------------------------
\subsection{Physical interpretation}
\label{sec:interpretation}

Equation~(\ref{eq:conjugate_dynamics}) has a simple physical meaning: the real part of the correlation matrix changes at a rate determined by the imaginary part, and vice versa. In the same way that position and momentum in classical mechanics exchange their roles under a Hamiltonian flow~--- $\dot{q} = \partial H/\partial p$, $\dot{p} = -\partial H/\partial q$~--- the two components of the quantum correlation matrix drive each other's evolution.

The ``gap'' between the Schr\"odinger equation and the Born rule is therefore not an empty conceptual void. It has internal dynamical structure: the real and imaginary parts of the correlation matrix are locked in a perpetual exchange mediated by $H$. The Born rule, which projects $C$ onto $C^R$ via $|\cdot|^2$, discards $C^I$ and thereby breaks this conjugate cycle. What is lost is not a passive phase variable but an active dynamical degree of freedom.

%% -----------------------------------------------------------------
\subsection{Relation to existing results}
\label{sec:related_existing}

The algebraic separation of real and imaginary parts of the Heisenberg equation is mathematically equivalent to the standard Bloch equation form for the single-particle density matrix~\cite{peschel2009}. The contribution of the present work is threefold: (a)~the recognition of the symplectic structure $\Omega = \sum dC^R \wedge dC^I$ and its conservation under the Heisenberg flow, (b)~the identification of $(C^R, C^I)$ as canonical conjugate pairs at the level of correlation matrix elements, and (c)~the use of this structure as the foundation for the dynamical necessity of complex numbers (Theorem~1, Sec.~\ref{sec:theorem1}). Items~(a) and~(b) constitute the new interpretative work; the algebraic form on which they rest is standard.

%% =================================================================
%% Section III: Theorem 1
%% =================================================================
\section{Theorem 1: Mode Decomposition Necessity of Complex Numbers}
\label{sec:theorem1}

%% -----------------------------------------------------------------
\subsection{Describability versus decomposability}
\label{sec:th1_describability}

Consider the linear dynamical system on $\mathbb{R}^2$:
\begin{equation}
\frac{d}{dt}\begin{pmatrix} A \\ B \end{pmatrix}
= \begin{pmatrix} 0 & \omega \\ -\omega & 0 \end{pmatrix}
\begin{pmatrix} A \\ B \end{pmatrix},
\qquad \omega \in \mathbb{R},\; \omega \neq 0. \label{eq:M_system}
\end{equation}
The general solution over $\mathbb{R}$ is
\begin{equation}
\begin{pmatrix} A(t) \\ B(t) \end{pmatrix} =
\begin{pmatrix} A_0\cos\omega t + B_0\sin\omega t \\
-A_0\sin\omega t + B_0\cos\omega t \end{pmatrix}. \label{eq:real_solution}
\end{equation}
Every quantity in Eq.~\eqref{eq:real_solution} is real: $A$, $B$, $A_0$, $B_0$, $\cos$, $\sin$. The system is describable over $\mathbb{R}$.

\textbf{Definition 1 (Describability over $\mathbb{R}$).} A linear dynamical system on $\mathbb{R}^n$ is describable over $\mathbb{R}$ if its time-dependent solutions can be written using only real-valued functions and real initial conditions.

The existence of real solutions is uncontroversial and well-known. The relevant question is different.

\textbf{Definition 2 (Decomposability over $\mathbb{R}$).} A linear dynamical system $d\mathbf{x}/dt = M\mathbf{x}$ on $\mathbb{R}^n$ is decomposable over $\mathbb{R}$ if there exists an invertible real matrix $P$ such that
\begin{equation}
P^{-1}MP = \operatorname{diag}(\lambda_1, \lambda_2, \dots, \lambda_n),\qquad \lambda_i \in \mathbb{R}. \label{eq:diag_real}
\end{equation}
In the transformed coordinates $\mathbf{u} = P^{-1}\mathbf{x}$, the dynamics decouples into $n$ independent one-dimensional equations $du_i/dt = \lambda_i u_i$. Each $u_i$ evolves autonomously, depending only on itself.

\textbf{Claim 1.} The system~\eqref{eq:M_system} is \textbf{not} decomposable over $\mathbb{R}$.

\textbf{Proof.} The characteristic polynomial of $M$ is $\det(M-\lambda I) = \lambda^2 + \omega^2 = 0$, giving eigenvalues $\lambda = \pm i\omega$. These are not elements of $\mathbb{R}$ for $\omega \neq 0$. A real matrix is diagonalizable over $\mathbb{R}$ iff all its eigenvalues are real. Since the eigenvalues are not real, no real $P$ satisfying~\eqref{eq:diag_real} exists. $\square$

The distinction is now formally established. Describability asks whether real-valued functions can parametrize the time evolution; decomposability asks whether the state space can be split into independently evolving one-dimensional invariant subspaces labeled by real eigenvalues. The system admits the first but not the second.

\textbf{Why decomposability matters for quantum mechanics.} Quantum mechanics is constitutively a spectral theory. Unlike classical mechanics, which can remain in the time domain and describe trajectories, quantum mechanics is built on:

\begin{enumerate}
\item[\textbf{1.}] \textbf{Energy eigenvalues.} The spectrum of $H$ defines the possible energies. Every spectral decomposition $H = \sum_n E_n |E_n\rangle\langle E_n|$ requires diagonalization.
\item[\textbf{2.}] \textbf{Occupation numbers.} In second quantization, $\hat{n}_k = c^\dagger_k c_k$ is diagonal in the single-particle eigenbasis. The occupation number operator is defined in the diagonal basis.
\item[\textbf{3.}] \textbf{Fock space.} The many-body Hilbert space is $\mathcal{F} = \bigoplus_N \mathcal{H}_N$, spanned by states $|n_1, n_2, \ldots\rangle$ with $\sum n_i = N$. Each occupation string consists of eigenvalues of the number operators in the diagonal basis.
\item[\textbf{4.}] \textbf{Quantum numbers.} Every quantum number (spin projection, angular momentum, momentum mode, phonon branch) labels a one-dimensional invariant subspace defined by diagonalization of some Hermitian operator.
\item[\textbf{5.}] \textbf{Born rule.} $P(n) = |\langle n|\psi\rangle|^2$ projects the state onto basis states $|n\rangle$ that are eigenstates of the number operator in a diagonal basis.
\end{enumerate}

In every case, the physical content of quantum mechanics is carried not by trajectory functions $A(t)$, $B(t)$ but by the labels of diagonalized subspaces---the eigenvalues $\lambda_i$ and the occupation numbers $n_k$. These require decomposability, not mere describability.

%% -----------------------------------------------------------------
\subsection{Lemma 1: Mode decomposition impossibility over $\mathbb{R}$}
\label{sec:th1_lemma1}

\textbf{Lemma 1 (Mode Decomposition Impossibility).} Let a linear dynamical system on $\mathbb{R}^2$ be governed by $M = \bigl(\begin{smallmatrix}0 & \omega \\ -\omega & 0\end{smallmatrix}\bigr)$ with $\omega \ne 0$. Then:

\textbf{(i)} There exists no real invertible matrix $P$ such that $P^{-1}MP$ is diagonal over $\mathbb{R}$.

\textbf{(ii)} There exists no nontrivial real linear functional $L(A,B) = \alpha A + \beta B$ that is invariant under the flow (i.e., $dL/dt = 0$).

\textbf{(iii)} There exists no decomposition of the two-dimensional real phase space into one-dimensional invariant subspaces over $\mathbb{R}$ such that the dynamics respects the decomposition.

\textbf{Proof.}

(i) The characteristic polynomial $\det(M - \lambda I) = \lambda^2 + \omega^2$ has discriminant $\Delta = -4\omega^2 < 0$. The eigenvalues $\lambda = \pm i\omega$ are complex for $\omega \neq 0$. A real matrix is diagonalizable over $\mathbb{R}$ iff all eigenvalues are real and the minimal polynomial splits over $\mathbb{R}$ with distinct roots. Since the eigenvalues are not real, diagonalization over $\mathbb{R}$ is impossible.

(ii) Let $L(A,B) = \alpha A + \beta B$ with $\alpha,\beta\in\mathbb{R}$. Then
\begin{equation}
\frac{dL}{dt} = \alpha\frac{dA}{dt} + \beta\frac{dB}{dt}
= \alpha\omega B + \beta(-\omega A)
= \omega(\alpha B - \beta A). \label{eq:Ldot}
\end{equation}
For $dL/dt = 0$ for all $(A,B)$, we require $\alpha = 0$ and $\beta = 0$. Hence only the trivial functional is invariant. Every nonzero linear functional varies under the flow---there is no conserved linear quantity.

(iii) Follows from (i): if a decomposition into one-dimensional invariant subspaces existed, each subspace would be spanned by a real eigenvector, and concatenating these eigenvectors would yield a diagonalizing matrix $P$. No real eigenvectors exist (the eigenvector equation $Mv = \pm i\omega v$ has no nonzero real solution). Therefore no such decomposition exists. $\square$

\textbf{Corollary 1 (Quantum Number Impossibility).} Since quantum numbers (occupation numbers $n_k$, phonon mode indices, angular momentum projections $m$, spin projections $s_z$) are eigenvalues of Hermitian operators that label one-dimensional invariant subspaces of the dynamics, they cannot be consistently defined within a purely real-valued description of canonically conjugate dynamics that does not admit an eigenbasis decomposition.

\textbf{Proof.} Each quantum number $n_k$ labels a one-dimensional subspace of the single-particle Hilbert space that is invariant under the free Hamiltonian $H_0$. This subspace is an eigenspace of the number operator $\hat{n}_k = c^\dagger_k c_k$, and its existence requires diagonalizing the single-particle dynamics. By Lemma~1(iii), this diagonalization is impossible over $\mathbb{R}$ for the canonical conjugate pair $(C^R_k,\,C^I_k)$. To assign a quantum number $n_k$ is to assert that the corresponding mode evolves independently---precisely the decomposability that $\mathbb{R}$ forbids.

The experimental fact that quantum systems exhibit discrete, well-defined quantum numbers---occupation numbers in quantum gas microscopes, Landau levels in the quantum Hall effect, phonon branches in neutron scattering---therefore constitutes an empirical falsification of $\mathbb{R}$-only quantum mechanics. Every measured quantum number is evidence that the underlying dynamics decomposes into independent modes, which requires the spectral decomposition that only $\mathbb{C}$ provides. $\square$

%% -----------------------------------------------------------------
\subsection{Why the harmonic oscillator confirms our framework}
\label{sec:th1_harmonic}

The objection that ``the classical harmonic oscillator has the same matrix $M$ and works over $\mathbb{R}$'' is not a counterexample---it is supporting evidence.

\paragraph{Classical physics in the time domain.} Classical mechanics can describe the harmonic oscillator trajectory using $\mathbb{R}$-valued functions because it is satisfied with trajectories. The description $x(t) = x_0\cos\omega t + (p_0/m\omega)\sin\omega t$ is a complete characterization. No mode decomposition is required because the classical state is fully specified by $(x,p)$ at each instant.

\paragraph{Classical physics itself flees to $\mathbb{C}$ when mode decomposition is needed.} The moment classical physics requires spectral decomposition---decomposing motion into independent frequency components---it abandons $\mathbb{R}$ and adopts $\mathbb{C}$. Three canonical examples:

\textbf{Example 1: Fourier analysis.} The Fourier transform of a real signal $f(t)$ is complex:
\begin{equation}
\tilde f(\omega) = \int_{-\infty}^{\infty} f(t)\,e^{-i\omega t}\,dt.
\end{equation}
The complex exponential basis $e^{-i\omega t}$ is used because it diagonalizes the time-translation operator. The real $\sin$ and $\cos$ basis does diagonalize the second-derivative operator, but the pair $(\sin,\cos)$ are coupled: $d/dt\sin = \cos$, $d/dt\cos = -\sin$---precisely the same $M$ matrix. Only the complex exponential $e^{-i\omega t}$ satisfies $d/dt(e^{-i\omega t}) = -i\omega(e^{-i\omega t})$, a one-dimensional eigen-equation.

\textbf{Example 2: Electrical impedance.} The impedance of a capacitor is $Z_C = 1/(i\omega C)$, of an inductor is $Z_L = i\omega L$. These are complex numbers. Real-valued circuit analysis (instantaneous voltage and current) is possible, but the moment one needs frequency response, transfer functions, or resonance---i.e., mode decomposition---complex impedance is mandatory.

\textbf{Example 3: Normal mode analysis.} For $N$ coupled oscillators, the normal mode frequencies $\omega_k$ are eigenvalues of the dynamical matrix. The eigenvectors are generally complex (they contain phase relationships between oscillators). While normal coordinates can be chosen real for undamped systems, the eigenvalue problem that finds the modes is necessarily over $\mathbb{C}$.

\paragraph{Quantum mechanics is constitutively spectral.} Quantum mechanics differs from classical mechanics in a fundamental way: it has no trajectory description to fall back on. The wavefunction $|\psi(t)\rangle$ is not a trajectory in configuration space---it is a vector in Hilbert space whose physical content is extracted through spectral decomposition. Measurement outcomes are eigenvalues (spectral theorem); probabilities are $|\langle\text{eigenstate}|\psi\rangle|^2$ (Born rule); time evolution is $e^{-iHt/\hbar}$ (spectral representation of the propagator); uncertainty relations are rooted in commutators $[X,P] = i\hbar$ (a complex structure).

\begin{table}[t]
\caption{Classical versus quantum harmonic oscillator. The classical oscillator avoids $\mathbb{C}$ by staying in the time domain; the quantum oscillator cannot.}
\label{tab:HO_comparison}
\begin{ruledtabular}
\begin{tabular}{lcc}
& Classical HO & Quantum HO \\
\hline
State description & $(x,p)$ trajectories & $|\psi\rangle = \sum c_n |n\rangle$ \\
$\mathbb{R}$-describable? & Yes ($\sin/\cos$) & Yes \\
$\mathbb{R}$-decomposable? & Not needed (trajectory suffices) & No---and this is fatal \\
Mode labels needed? & No (single particle) & Yes (Fock states, energy levels) \\
$\mathbb{C}$ used when? & Convenience (impedance, Fourier) & Necessity (spectral theorem) \\
\end{tabular}
\end{ruledtabular}
\end{table}

%% -----------------------------------------------------------------
\subsection{Exclusion of competing algebras under unitarity}
\label{sec:th1_algebra}

A mathematically sophisticated objection might be: ``There are three two-dimensional $\mathbb{R}$-algebras: $\mathbb{C}\;(i^2 = -1)$, split-complex $(j^2 = +1)$, and dual numbers $(\epsilon^2 = 0)$. Why is $\mathbb{C}$ uniquely selected?''

The answer requires imposing the physical constraint of unitarity (probability conservation).

\paragraph{The three two-dimensional $\mathbb{R}$-algebras.} For the evolution matrix $M = \bigl(\begin{smallmatrix}0 & \omega \\ -\omega & 0\end{smallmatrix}\bigr)$, consider the extension of $\mathbb{R}$ to each two-dimensional $\mathbb{R}$-algebra:

\textbf{Algebra A1: $\mathbb{C}$ ($i^2 = -1$).} Define $Z = C^R + iC^I$. Then $dZ/dt = -i\omega Z$, giving $Z(t) = Z(0)\,e^{-i\omega t}$. The norm $|Z|^2 = (C^R)^2 + (C^I)^2$ is constant. \textbf{Unitary. Probability-conserving.}

\textbf{Algebra A2: Split-complex ($j^2 = +1$).} Define $Z = C^R + jC^I$. The ``rotation'' generated by $M$ is hyperbolic: $dZ/dt = j\omega Z$, giving $Z(t) = Z(0)\,e^{j\omega t} = Z(0)[\cosh\omega t + j\sinh\omega t]$. The ``norm'' in the split-complex algebra is $|Z|_s^2 = (C^R)^2 - (C^I)^2$, which evolves as $d|Z|_s^2/dt = 2\omega|Z|_s^2 \neq 0$. \textbf{Norm is not conserved.} The hyperbolic rotation produces exponential growth or decay---a dynamical instability that violates unitarity and probability conservation.

\textbf{Algebra A3: Dual numbers ($\epsilon^2 = 0$).} Define $Z = C^R + \epsilon C^I$. The matrix $M$ in the dual-number representation is nilpotent: $d^2Z/dt^2 = 0$, giving $Z(t) = Z(0) + t\cdot\dot Z(0)$. The trajectory is linear in $t$---unbounded growth. Probability is not conserved.

\paragraph{Lemma 2 (Unitarity Constraint).} Let a linear dynamical system on $\mathbb{R}^2$ be governed by $M = \bigl(\begin{smallmatrix}0 & \omega \\ -\omega & 0\end{smallmatrix}\bigr)$ with $\omega \neq 0$, and let the evolution be required to conserve the Euclidean norm $\|(A,B)\|^2 = A^2 + B^2$ (the physical requirement of probability conservation). Then the unique two-dimensional $\mathbb{R}$-algebra that simultaneously (i) diagonalizes the dynamics (i.e., makes each mode evolve independently) and (ii) conserves the norm (i.e., generates unitary evolution) is $\mathbb{C}$.

\textbf{Proof.} The three two-dimensional $\mathbb{R}$-algebras $(\mathbb{C}$, split-complex, dual$)$ are the only finite-dimensional associative division algebras over $\mathbb{R}$ of dimension~2 (up to isomorphism). For each:

\begin{itemize}
\item $\mathbb{C}$: Eigenvalues $\pm i\omega$ generate $\mathrm{SO}(2)$ rotations, preserving $A^2 + B^2$. Unitary. Diagonalizing.
\item Split-complex: Eigenvalues $\pm\omega$ (real) generate hyperbolic rotations, preserving $A^2 - B^2$ but not $A^2 + B^2$. Non-unitary.
\item Dual: No eigenvalues (nilpotent Jordan block). Produces secular growth. Non-unitary, non-diagonalizing.
\end{itemize}

Only $\mathbb{C}$ satisfies both (i) and (ii). $\square$

Note that this argument does not assume $\mathbb{C}$ a priori. It examines all possible two-dimensional $\mathbb{R}$-algebra extensions and eliminates all but $\mathbb{C}$ through the unitarity constraint. The physical origin of the unitarity constraint is the conservation of total probability $\sum_n |c_n|^2 = 1$, which is equivalent to the conservation of $\operatorname{Tr}\rho = 1$ under the dynamics $dC/dt = i[H,C]$.

%% -----------------------------------------------------------------
\subsection{Lemma 3: Spectral Correspondence Between $M$ and $H$}
\label{sec:th1_lemma3}

\textbf{Lemma 3 (Spectral Correspondence).} Let $H = \sum_{k,l} h_{kl} c^\dagger_k c_l$ be a quadratic Hamiltonian with $h$ a real symmetric $N \times N$ matrix with eigenvalues $\{E_k\}_{k=1}^N$ and eigenvectors $\{|k\rangle\}_{k=1}^N$. Let $M_{\text{total}} = \bigoplus_{k<l} M_{kl}$ be the generator of the $(C^R, C^I)$ dynamics, expressed in the eigenbasis of $h$. Then:

\textbf{(i)} The eigenvalues of $M_{\text{total}}$ are $\{\pm i\omega_{kl}\}_{k<l}$ where $\omega_{kl} = (E_k - E_l)/\hbar$.

\textbf{(ii)} Each eigenvector of $M_{\text{total}}$ over $\mathbb{C}$ corresponds one-to-one with the single-particle excitation operator $|l\rangle\langle k| \equiv c^\dagger_l c_k$.

\textbf{(iii)} The spectral decomposition of $M_{\text{total}}$ over $\mathbb{C}$ is equivalent to the identification of all independent single-particle excitation modes of the system.

\paragraph{Proof of (i).} The block $M_{kl}$ has characteristic polynomial $\det(M_{kl} - \lambda I) = \lambda^2 + \omega_{kl}^2 = 0$, giving eigenvalues $\lambda = \pm i\omega_{kl}$. Since $M_{\text{total}}$ is block-diagonal, its eigenvalues are the union of the eigenvalues of each block:
\[
\operatorname{Spec}(M_{\text{total}}) = \{\pm i\omega_{kl} : 1 \leq k < l \leq N\}.
\]
The eigenfrequencies $\omega_{kl} = (E_k - E_l)/\hbar$ are the Bohr frequencies of $H$ --- the physically measurable energy differences between single-particle levels.

\paragraph{Proof of (ii).} Consider the complexification $Z_{kl} \equiv C^R_{kl} + i C^I_{kl} = \langle c^\dagger_l c_k \rangle$. From the dynamics,
\begin{equation}
\frac{d}{dt} Z_{kl} = i\omega_{kl} Z_{kl} \quad\Longrightarrow\quad Z_{kl}(t) = Z_{kl}(0) e^{i\omega_{kl} t},
\label{eq:lemma3_Zevol}
\end{equation}
while the operator $c^\dagger_l c_k = |l\rangle\langle k|$ satisfies the same Heisenberg equation:
\begin{equation}
\frac{d}{dt} (c^\dagger_l c_k) = \frac{i}{\hbar}[H, c^\dagger_l c_k] = i\omega_{kl} \, c^\dagger_l c_k.
\label{eq:lemma3_opevol}
\end{equation}
Hence $Z_{kl}$ is precisely the expectation value of the single-particle excitation operator $|l\rangle\langle k|$, and the eigenvector $(1,i)^\top$ of $M_{kl}$ (eigenvalue $-i\omega_{kl}$) represents this same physical mode. Each eigenvalue $\pm i\omega_{kl}$ therefore labels a distinct single-particle excitation.

\paragraph{Proof of (iii).} The set $\{Z_{kl}\}_{k<l}$ spans a $D$-dimensional complex vector space where $D = N(N-1)/2$. The dynamics in this space is fully diagonalized: $dZ_{kl}/dt = i\omega_{kl} Z_{kl}$ with no coupling between different $(k,l)$ pairs. The spectral decomposition of $M_{\text{total}}$ over $\mathbb{C}$ therefore yields precisely the set of independent single-particle excitation modes $|l\rangle\langle k|$. Conversely, any attempt to identify independent modes over $\mathbb{R}$ would require real eigenvectors of $M_{\text{total}}$; by Lemma~1(i), no such real eigenvectors exist. $\square$

%% -----------------------------------------------------------------
\subsection{Proposition 1: No Real Independent Excitation Modes}
\label{sec:th1_prop1}

\textbf{Proposition 1.} Let $M_{kl}$ be the generator of correlation dynamics for the pair $(k,l)$ with $\omega_{kl} \neq 0$. Then the following are equivalent:

\textbf{(a)} $M_{kl}$ has no real eigenvectors (i.e., is not diagonalizable over $\mathbb{R}$).

\textbf{(b)} There exists no nontrivial real linear functional $L(C^R_{kl}, C^I_{kl}) = \alpha C^R_{kl} + \beta C^I_{kl}$ satisfying $dL/dt = \lambda L$ for any $\lambda \in \mathbb{R}$.

\textbf{(c)} The single-particle excitation $|l\rangle\langle k| = c^\dagger_l c_k$ has no independent representation within the real vector space $\mathbb{R}^2$ spanned by $(C^R_{kl}, C^I_{kl})$.

\paragraph{Proof that (a) $\Leftrightarrow$ (b).} The eigenvalues of $M_{kl}$ are $\pm i\omega_{kl} \notin \mathbb{R}$. A real eigenvector requires a real eigenvalue; since none exist, statement~(a) holds. From Lemma~1(ii), for any $(\alpha,\beta) \neq (0,0)$, $dL/dt = \omega(\alpha C^I_{kl} - \beta C^R_{kl})$ cannot equal $\lambda L$ for any $\lambda \in \mathbb{R}$, hence (b) holds. The converse follows similarly.

\paragraph{Proof that (b) $\Leftrightarrow$ (c).} The operator $c^\dagger_l c_k$ creates a single-particle excitation from level $k$ to level $l$. Its expectation value evolves as $d\langle c^\dagger_l c_k\rangle/dt = i\omega_{kl} \langle c^\dagger_l c_k\rangle$, a one-dimensional complex evolution. If a real linear combination $L$ existed that evolved independently ($dL/dt = \lambda L$, $\lambda \in \mathbb{R}$), then $L$ would serve as a real representation of the excitation. Statement~(b) denies the existence of such $L$, hence no real representation exists within $(C^R_{kl}, C^I_{kl})$ space. $\square$

\paragraph{Interpretation.} Proposition~1 establishes the first link in the bridge: the failure of $M$ to be diagonalizable over $\mathbb{R}$ is equivalent to the absence of a real representation for any single-particle excitation. The excitation operator $c^\dagger_l c_k$, which carries the physical content of the quantum transition between levels $k$ and $l$, has no real-valued counterpart in the correlation dynamics.

%% -----------------------------------------------------------------
\subsection{Proposition 2: Occupation Numbers Require the Complex Mode Basis}
\label{sec:th1_prop2}

\textbf{Proposition 2.} Consider the quadratic Hamiltonian $H = \sum_{kl} h_{kl} c^\dagger_k c_l$ with single-particle energies $\{E_k\}$ and eigenbasis $\{|k\rangle\}$. The occupation number operator $\hat{n}_k = c^\dagger_k c_k$ and its eigenvalues $n_k \in \{0,1\}$ label the independent single-particle modes of the system. Under the dynamics generated by $H$, these mode labels can be consistently defined only if the dynamical independence of each mode is established. By Proposition~1, mode $k$ is dynamically independent from mode $l$ iff the coherence $C_{kl}$ can be represented as an independently evolving quantity, which requires $\mathbb{C}$.

\paragraph{Proof.} The occupation $\langle \hat{n}_k \rangle = C_{kk}$ evolves as
\begin{equation}
\frac{d}{dt} C_{kk} = \frac{i}{\hbar}\sum_{l \neq k} \bigl( h_{lk} C_{lk} - h_{kl} C_{kl} \bigr).
\label{eq:prop2_occ_evol}
\end{equation}
Each term $C_{kl}$ in this sum is a coherence involving mode $k$. For mode $k$ to be an identifiable independent degree of freedom, each $C_{kl}$ must be decomposable into an independent dynamical channel. In the real representation, $(C^R_{kl}, C^I_{kl})$ are permanently coupled by $M_{kl}$ --- neither evolves independently, and the sum in Eq.~\eqref{eq:prop2_occ_evol} involves variables whose collective dynamics cannot be decomposed into independent mode contributions. Over $\mathbb{C}$, $Z_{kl} = C^R_{kl} + i C^I_{kl}$ evolves as an independent one-dimensional system $dZ_{kl}/dt = i\omega_{kl} Z_{kl}$, and the contribution of the $(k,l)$ coherence to $C_{kk}$ becomes an independent term. Hence the mode label $k$ is meaningful only when the complex representation is available. $\square$

\paragraph{Corollary 2.} The entire Fock space construction $\mathcal{F} = \bigoplus_N \mathcal{H}_N$, whose basis states $|n_1, n_2, \ldots\rangle$ are labeled by occupation numbers, rests on the prior identification of independent single-particle modes $\{|k\rangle\}$. Since that identification requires $\mathbb{C}$ (Proposition~2), the Fock space itself cannot be consistently defined within a purely real-valued description.

%% -----------------------------------------------------------------
\subsection{Bridge Theorem: The Complete Equivalence}
\label{sec:th1_bridge}

We now assemble the chain of implications into a single theorem that directly addresses the connection between the algebraic properties of $M$ and the necessity of $\mathbb{C}$ for quantum mechanics.

\textbf{Bridge Theorem.} For a quadratic fermionic Hamiltonian $H = \sum_{kl} h_{kl} c^\dagger_k c_l$ with non-degenerate spectrum, let $M_{\text{total}} = \bigoplus_{k<l} M_{kl}$ be the generator of correlation dynamics. Then the following are equivalent:

\begin{enumerate}
\item[\textbf{(i)}] $M_{\text{total}}$ is diagonalizable over $\mathbb{R}$.
\item[\textbf{(ii)}] Every single-particle excitation operator $c^\dagger_l c_k$ can be represented as an independently evolving real quantity.
\item[\textbf{(iii)}] The occupation number operators $\hat{n}_k$ can be defined as dynamically independent labels within a purely real framework.
\item[\textbf{(iv)}] The Fock space $\mathcal{F}$ can be constructed without leaving $\mathbb{R}$.
\item[\textbf{(v)}] Quantum mechanics can be formulated over $\mathbb{R}$.
\end{enumerate}

\paragraph{Proof.} The equivalence chain is:
\begin{align*}
\text{(i) fails} &\Longleftrightarrow \text{Eigenvalues of } M_{\text{total}} \text{ are not all real} \quad (\text{Lemma~1(i)})\\
\text{(ii) fails} &\Longleftrightarrow \text{No real } L \text{ satisfies } dL/dt = \lambda L,\; \lambda\in\mathbb{R} \quad (\text{Proposition~1})\\
\text{(iii) fails} &\Longleftrightarrow \text{Mode labels } n_k \text{ have no independent dynamics} \quad (\text{Proposition~2})\\
\text{(iv) fails} &\Longleftrightarrow \text{Fock basis cannot be consistently defined} \quad (\text{Corollary~2})\\
\text{(v) fails} &\Longleftrightarrow \text{Real QM is inconsistent} \quad (\text{QM = Fock space + dynamics})
\end{align*}
Since (i) is false --- the eigenvalues of $M_{\text{total}}$ are $\pm i\omega_{kl} \notin \mathbb{R}$ for any $\omega_{kl} \neq 0$ --- it follows that (v) is false: quantum mechanics cannot be formulated over $\mathbb{R}$. Conversely, if $M_{\text{total}}$ were diagonalizable over $\mathbb{R}$, each $(C^R_{kl}, C^I_{kl})$ pair would decouple into real independent modes, enabling the entire construction within $\mathbb{R}$. $\square$

\paragraph{Direct address of the category error.}
A natural objection is that $M$ and $H$ are different operators acting on different spaces, so the non-diagonalizability of $M$ cannot imply anything about $H$'s spectral structure. This objection rests on a misunderstanding of the relationship between $M$ and $H$, which Lemma~3 resolves. $M$ is not claimed to equal $H$; rather, $M$ is the \textbf{real representation of the adjoint action} $\operatorname{ad}_H = [H,\cdot]$ restricted to the single-particle correlation sector. The relationship is:

\begin{equation}
H \xrightarrow{\text{Heisenberg eq.}} \frac{dC}{dt} = \frac{i}{\hbar}[H,C] \xrightarrow{\text{real/imag. decomp.}} \frac{d}{dt}\binom{C^R}{C^I} = M\binom{C^R}{C^I}.
\label{eq:bridge_chain}
\end{equation}

$M$ is the push-forward of $H$ through the Heisenberg equation and the real/imaginary decomposition. Its eigenvalues are the Bohr frequencies $\pm i\omega_{kl} = \pm i(E_k - E_l)/\hbar$ --- the energy differences of $H$. The spectral correspondence (Lemma~3) establishes a bijection between the complex eigenmodes of $M$ and the single-particle excitations of $H$. Table~\ref{tab:M_identity} clarifies what $M$ is and is not.

\begin{table}[t]
\caption{What $M$ IS and IS NOT.}
\label{tab:M_identity}
\begin{ruledtabular}
\begin{tabular}{lp{0.65\linewidth}}
Statement & Truth value \\
\hline
$M = H$ & \textbf{False} --- $M$ is $2N(N-1)\times 2N(N-1)$; $h$ is $N\times N$ \\
$M$ acts on the same space as $H$ & \textbf{False} --- $M$ acts on $(C^R,C^I)\in\mathbb{R}^{N(N-1)}$; $H$ acts on $\mathcal{H}$ \\
$M$ is unrelated to $H$ & \textbf{False} --- $M$ is the real representation of $\operatorname{ad}_H$ restricted to correlations \\
$M$'s eigenvalues are $H$'s energy differences & \textbf{True} --- $\operatorname{Spec}(M) = \{\pm i(E_k-E_l)/\hbar\}_{k<l}$ \\
$M$'s complex eigenvectors ARE the excitation operators & \textbf{True} --- $Z_{kl} = C^R_{kl}+iC^I_{kl} = \langle c^\dagger_l c_k\rangle$ \\
$M$ not $\mathbb{R}$-diagonalizable $\Rightarrow$ excitations not separable over $\mathbb{R}$ & \textbf{True} --- Lemma~3 + Proposition~1 \\
Excitations not separable $\Rightarrow$ occupation numbers not definable over $\mathbb{R}$ & \textbf{True} --- Proposition~2 \\
Occupation numbers not definable $\Rightarrow$ QM requires $\mathbb{C}$ & \textbf{True} --- Bridge Theorem \\
\end{tabular}
\end{ruledtabular}
\end{table}

%% -----------------------------------------------------------------
\subsection{Relation to existing work}
\label{sec:th1_relation}

\textbf{Distinction from de Oliveira~(2023).} de Oliveira~\cite{deoliveira2023} showed that canonical variables in classical Hamiltonian mechanics can be complexified, yielding a representation that is algebraically convenient. The conclusion is that complex numbers are \textit{natural} for describing canonical dynamics. Our result is stronger: starting from the canonical conjugate structure of the correlation matrix---which is forced by the Heisenberg equation, not chosen for convenience---we prove that the real evolution matrix is not decomposable over $\mathbb{R}$ because its eigenvalues are imaginary. A real-valued description is not merely inconvenient but mathematically incomplete for the spectral decomposition that quantum mechanics requires. The distinction is between ``natural'' and ``necessary,'' and the necessity arises from the algebraic incompleteness of $\mathbb{R}$ under decomposability.

\textbf{Distinction from Renou~\textit{et al.}~(2021).} Renou~\textit{et al.}~\cite{renou2021} showed that real and complex quantum mechanics make operationally different predictions in network Bell scenarios. Their result answers the empirical question: ``Does real quantum mechanics reproduce all statistical predictions of complex quantum mechanics?'' The answer is no. Our result addresses the structural question: ``Why must real quantum mechanics fail?'' The answer is that the canonical conjugate dynamics of the correlation matrix cannot be decomposed over $\mathbb{R}$---the operational failure identified by Renou~\textit{et al.} has a structural root in the algebraic incompleteness of $\mathbb{R}$ as a field for decomposing correlation matrix dynamics into independent spectral modes. Where Renou answers ``does real QM fail?'' we answer ``why must real QM fail?'' The two results are complementary: one establishes the fact of operational inequivalence, the other traces it to a dynamical necessity.

\textbf{Implicit in the framework.} Theorem~1 does not assert that complex numbers are a priori necessary for all of quantum mechanics---such a claim would be both too broad and philosophically charged. The theorem asserts that, \textit{given} the experimentally observed complex-valued correlation matrix of fermionic systems and \textit{given} its Heisenberg dynamics (Identity~1), a purely real-valued description of this dynamics cannot achieve the spectral decomposition that quantum mechanics requires for occupation numbers, energy levels, and quantum numbers. The complex structure is not an axiom we impose; it is dynamically forced by the algebraic structure of the Heisenberg evolution.

%% =================================================================
%% Section IV: Theorem 2
%% =================================================================
\section{Theorem 2: Ergotropy Rate Decomposition}
\label{sec:theorem2}

%% -----------------------------------------------------------------
\subsection{Proposition and derivation}
\label{sec:th2_derivation}

\textbf{Proposition.} For any finite-dimensional quantum system, the ergotropy rate admits the exact decomposition

\begin{equation}
\boxed{\dot{\mathcal{E}} = -\sum_{m\neq n} \omega_{mn} (\Delta H_\rho)_{nm} C^I_{mn}},
\label{eq:ergo_decomp}
\end{equation}

where $\Delta H_\rho = H - H_\rho^{\text{pass}}$ is the difference between the Hamiltonian and the passive state Hamiltonian, $\omega_{mn} = (E_m - E_n)/\hbar$, and $C^I_{mn}$ is the imaginary part of the correlation matrix element.

\textbf{Derivation.} The ergotropy rate equation established by Huang~\cite{huang2026} is

\begin{equation}
\dot{\mathcal{E}} = \operatorname{Tr}[\Delta H_\rho \cdot \dot{\rho}].
\label{eq:ergo_rate_Huang}
\end{equation}

In the eigenbasis of $H$, the density matrix evolves as $\dot{\rho}_{mn} = i\omega_{mn} \rho_{mn}$. Expressing $\rho$ in terms of the correlation matrix~--- for fermionic systems the two are related by $\rho_{mn} = C_{mn}$ for $m \neq n$ (the diagonal elements are the occupation numbers $\langle n_m \rangle$)~--- gives

\begin{equation}
\dot{\mathcal{E}} = \sum_{m,n} (\Delta H_\rho)_{nm} \dot{\rho}_{mn}
= \sum_{m\neq n} (\Delta H_\rho)_{nm} \cdot i\omega_{mn}(C^R_{mn} + i C^I_{mn}).
\label{eq:ergo_expand}
\end{equation}

Taking the real part (the ergotropy rate is a physical observable) yields Eq.~(\ref{eq:ergo_decomp}). The passive state Hamiltonian $H_\rho^{\text{pass}} = \sum_k \epsilon_{\sigma(k)} |r_k\rangle\langle r_k|$ depends only on the eigenvalues $\{|a_n|^2\} = \{C^R_{nn}\}$ of $\rho$ and not on the phases, so $\Delta H_\rho$ is determined entirely by $C^R$.

%% -----------------------------------------------------------------
\subsection{Physical interpretation}
\label{sec:th2_interpretation}

Equation~(\ref{eq:ergo_decomp}) expresses a structural relationship: the rate of change of extractable work is driven exclusively by the imaginary part of the correlation matrix. The real part $C^R$ determines the ``capacity'' (through $\Delta H_\rho$), while the imaginary part $C^I$ provides the ``current.'' A purely mixed state ($C^I = 0$) has zero ergotropy current regardless of its energy content; a coherent state ($C^I \neq 0$) can deliver power even at the same energy expectation value.

This decomposition carries a striking implication. The Born rule $|\cdot|^2 : \mathbb{C} \to \mathbb{R}_{\ge 0}$ projects the correlation matrix onto its real-valued occupation spectrum, discarding $C^I$. Theorem~2 shows that what is discarded is precisely the quantity that drives extractable work. The Born--Schr\"odinger gap is therefore not a defect of the quantum formalism~--- it is a necessary structural feature that enables quantum thermodynamic advantage. The gap is the operational source of quantum power.

%% -----------------------------------------------------------------
\subsection{Operational protocol for $C^I$ measurement}
\label{sec:th2_protocol}

Equation~(\ref{eq:ergo_decomp}) connects $\dot{\mathcal{E}}$ to $C^I$, but if $C^I$ is only known through the density matrix $\rho$, the relation is a mathematical identity rather than an independent physical law. To break this circularity, an operational protocol for measuring $C^I$ without passing through the Born rule projection on the system is required.

A weak measurement protocol based on the weak-value formalism~\cite{dressel2014} provides such a path. In this protocol, $C^I_{mn}$ is extracted from the phase shift of an ancilla qubit coupled weakly to the system operator $\hat{B}_{mn} = i(c^\dagger_n c_m - c^\dagger_m c_n)$. The pointer readout uses projective measurement on the ancilla (which is unavoidable), but the system state is minimally disturbed, and the ancilla shift is proportional to $\langle \hat{B}_{mn} \rangle = 2i C^I_{mn}$ before any system projection occurs. The protocol therefore provides an operationally distinct route to $C^I$ that does not go through $|\langle c^\dagger_n c_m \rangle|^2$.

The full protocol, together with platform-specific implementations (superconducting qubits, cold-atom quantum gas microscopes), experimental feasibility estimates, and the two-copy verification strategy that elevates Eq.~(\ref{eq:ergo_decomp}) from identity to testable physical law, is presented in Sec.~\ref{sec:weakmeasurement}.

%% =================================================================
%% Section V: Weak Measurement Protocol
%% =================================================================
\section{Operational Independence: Weak Measurement Protocol for $C^I$}
\label{sec:weakmeasurement}

%% -----------------------------------------------------------------
\subsection{The circularity and its resolution}
\label{sec:circularity}

Equation~(\ref{eq:ergo_decomp}) expresses the ergotropy rate $\dot{\mathcal{E}}$ as a linear combination of imaginary correlation matrix elements $C^I_{mn}$. If $C^I_{mn}$ is known only through the density matrix $\rho$~--- via $C^I_{mn} = \operatorname{Im}(\operatorname{Tr}[\rho c^\dagger_n c_m])$~--- then Eq.~(\ref{eq:ergo_decomp}) is a mathematical identity following from the definitions of ergotropy and the correlation matrix. It becomes a falsifiable physical law only when $C^I$ and $\dot{\mathcal{E}}$ are accessed through operationally independent measurement chains.

We emphasize at the outset that no measurement protocol can fully ``escape'' the Born rule in an absolute sense. The ancilla readout in the protocol below is itself projective, and any expectation value is ultimately a statistical average of individual measurement outcomes. The operational distinction we draw is more subtle. In the standard approach, $C^I_{mn} = \operatorname{Im}(\operatorname{Tr}[\rho c^\dagger_n c_m])$ is computed \textit{from} $\rho$, which is itself reconstructed via quantum state tomography~--- a process that applies the Born rule to the system's degrees of freedom at every measurement step. In the weak measurement protocol, the ancilla shift $\langle \sigma_y \rangle_{\text{anc}} \propto \langle \hat{B}_{mn} \rangle_{\text{sys}}$ is proportional to the system's \textit{operator expectation value}, not to Born probabilities of the system's observables. The system state $\rho$ is never reconstructed, and the system's degrees of freedom are never projectively measured. The ancilla plays the role of a ``witness'' that samples the pre-measurement expectation without forcing the system into an eigenstate.

We also note that the independent measurement of $\dot{\mathcal{E}}$ requires computing $\Delta H_\rho$, which in turn requires the eigenvalues of $\rho$. This step does use the Born rule. The operational independence is therefore between $[C^I$ measured via ancilla weak coupling$]$ and $[\dot{\mathcal{E}}$ measured via time-series tomography$]$~--- two measurement chains that share no intermediate steps. The theorem is falsifiable because an error in either chain would cause Eq.~(\ref{eq:ergo_decomp}) to fail.

%% -----------------------------------------------------------------
\subsection{Protocol design}
\label{sec:protocol_design}

\paragraph{Step 1: System-pointer coupling.}
Consider a fermionic system on $N$ lattice sites prepared in the state of interest. For a chosen pair of sites $(m,n)$, define the Hermitian operator

\begin{equation}
\hat{B}_{mn} = i(c^\dagger_n c_m - c^\dagger_m c_n),
\label{eq:B_operator}
\end{equation}

whose expectation value gives $\langle \hat{B}_{mn} \rangle = -2 C^I_{mn}$. Couple the system to an ancillary two-level system (pointer qubit) via the interaction Hamiltonian

\begin{equation}
H_{\text{int}} = \hbar g(t)\, \hat{B}_{mn} \otimes \hat{\sigma}_z^{\text{(anc)}},
\label{eq:int_Hamiltonian}
\end{equation}

where $g(t)$ is a controllable coupling strength satisfying $\int_0^\tau g(t)\,dt = g\tau \ll 1/\Delta \hat{B}_{mn}$ (the weak measurement condition). The interaction preserves the pointer's $\hat{\sigma}_z$ eigenbasis while the system operator $\hat{B}_{mn}$ generates a conditional phase shift.

\paragraph{Step 2: Weak measurement readout.}
Initialize the ancilla in the state $|+\rangle_x = (|0\rangle + |1\rangle)/\sqrt{2}$. After the interaction time $\tau$, the combined system-ancilla state is, to first order in $g\tau$,

\begin{equation}
|\Psi(\tau)\rangle \approx \bigl(1 - i g\tau \hat{B}_{mn} \otimes \hat{\sigma}_z\bigr) |\psi_{\text{sys}}\rangle \otimes |+\rangle_x.
\label{eq:weak_state}
\end{equation}

Measure the ancilla in the $\hat{\sigma}_y$ basis. The expectation value of $\hat{\sigma}_y$ on the ancilla is

\begin{equation}
\langle \hat{\sigma}_y \rangle_{\text{anc}} = 2g\tau \langle \hat{B}_{mn} \rangle_{\text{sys}} + O\bigl((g\tau)^2\bigr).
\label{eq:sigma_y}
\end{equation}

\paragraph{Step 3: Extract $C^I$.}
From the ancilla readout,

\begin{equation}
C^I_{mn} = -\frac{\langle \hat{\sigma}_y \rangle_{\text{anc}}}{4g\tau} + O(g\tau).
\label{eq:CI_extract}
\end{equation}

The weak measurement does not project the system onto an eigenstate of $\hat{B}_{mn}$. The ancilla shift is proportional to $\langle \hat{B}_{mn} \rangle$, which is the \textit{expectation value} carried by the pre-existing state, not an eigenvalue extracted from the system. The phase information is read from the ancilla without applying $|\cdot|^2$ to the system's degrees of freedom.

\paragraph{Step 4: Independent measurement of $\dot{\mathcal{E}}$.}
On an independent, identically prepared copy of the system (copy~B), perform full quantum state tomography at times $t$ and $t + \Delta t$. From $\rho(t)$, compute the passive state Hamiltonian $H_\rho^{\text{pass}} = \sum_k \epsilon_{\sigma(k)} |r_k\rangle\langle r_k|$ (which depends only on the eigenvalues $|a_n|^2 = C^R_{nn}$) and the ergotropy $\mathcal{E}(t) = \operatorname{Tr}[\rho(t)H] - \operatorname{Tr}[H_\rho^{\text{pass}}]$. Finite-differencing gives $\dot{\mathcal{E}} \approx [\mathcal{E}(t+\Delta t) - \mathcal{E}(t)]/\Delta t$.

\paragraph{Step 5: Falsification check.}
With $C^I_{mn}$ from Step~3 and $\dot{\mathcal{E}}$ from Step~4, compute

\begin{equation}
\dot{\mathcal{E}}_{\text{predicted}} = -\sum_{m\neq n} \omega_{mn} (\Delta H_\rho)_{nm} C^I_{mn}
\label{eq:predicted}
\end{equation}

and compare with $\dot{\mathcal{E}}_{\text{measured}}$. If they agree within experimental uncertainty, Theorem~2 is confirmed as a physical law relating two independently measured quantities. If they disagree beyond statistical error, Theorem~2 is falsified, and the theoretical framework must be revised.

%% -----------------------------------------------------------------
\subsection{Platform implementations}
\label{sec:platforms}

\textbf{Superconducting qubits.}
For a system of three transmon qubits implementing the single-spin-device engine (see the full text, Secs.~VI--IX), the weak measurement of $C^I_{12}$ (between the ground and first excited states) can be achieved via dispersive readout. A readout resonator coupled dispersively to the qubit pair provides cross-Kerr interaction $H_{\chi} = \hbar\chi_{12} \hat{n}_1 \hat{n}_2$. Driving the common resonator with a weak tone produces a reflected signal whose quadrature component is proportional to $\langle \hat{B}_{12} \rangle = -2 C^I_{12}$. Typical parameters: $g/2\pi \sim 1$ MHz, $\tau \sim 100$ ns, weak-measurement condition $g\tau \sim 0.1$. Single-shot SNR for $C^I$ is $\sim 0.3$, requiring $\sim 100$ repetitions for $3\%$ precision. Full tomography of the three-level system requires 8 independent measurements in the Gell-Mann basis, taking $\sim 1\,\mu$s per point. Total measurement time per parameter point is approximately 10 seconds, within reach of present-day superconducting qubit platforms.

\textbf{Cold-atom quantum gas microscopes.}
Fermionic quantum gas microscopes using $^6$Li atoms have demonstrated site-resolved density and spin correlation measurements~\cite{parsons2015,parsons2016,mazurenko2017}. The measurement of off-diagonal correlations $\langle c^\dagger_i c_j \rangle$ in fermionic systems has been achieved through lattice modulation spectroscopy~\cite{cocchi2017} and pair correlation analysis~\cite{yao2025}. For the weak measurement of $C^I$, a short $\pi/2$ rotation pulse maps the quadrature $\hat{B}_{mn}$ onto the density-measurable quadrature $\hat{A}_{mn} = c^\dagger_n c_m + c^\dagger_m c_n$, converting $C^I_{mn}$ into a density correlation $C^R_{mn}$ that can be read out via fluorescence imaging. The conversion factor is calibrated by the pulse area. Repeating the measurement without the $\pi/2$ pulse gives the original $C^R_{mn}$. While demonstrated for bosonic systems~\cite{karch2024}, the kinetic operator readout technique applies in principle to fermions when interactions are turned off during the measurement sequence. Site-resolved readout with $\sim 95\%$ imaging fidelity and $\sim 500$ shots yields $C^I$ precision of approximately $5\%$. The independent ergotropy measurement uses full counting statistics of site-resolved densities to reconstruct the single-particle density matrix and extract $\dot{\mathcal{E}}$.

%% -----------------------------------------------------------------
\subsection{Caveats and limitations}
\label{sec:caveats}

We note several practical considerations. First, the protocol requires site-resolved addressing of individual bilinear fermionic operators $c^\dagger_n c_m$, which remains technically challenging on current experimental platforms. For nearest-neighbor pairs in an optical lattice, this reduces to a measurement of the kinetic energy operator, which is accessible via lattice modulation spectroscopy~\cite{karch2024}. For longer-range pairs, the implementation is more demanding. Second, the weak measurement condition $g\tau \ll 1/\Delta\hat{B}_{mn}$ must be verified for each operator pair separately, as the operator variance $\Delta\hat{B}_{mn}$ depends on the state. Third, the need for independent, identically prepared copies assumes state-preparation fidelity that is high enough to ensure that the two copies are statistically indistinguishable~--- a standard assumption in quantum information experiments but one that must be validated in practice.

The protocol as presented is a proposal whose feasibility analysis is based on current technology trends. Actual implementation would benefit from collaboration with experimental groups specializing in weak measurement and quantum gas microscopy.

%% =================================================================
\section{Prediction A: State-Dependent Uncertainty Relation}
%% =================================================================

\subsection{Theorem and derivation}

Define the Hermitian operators
\begin{equation}
\hat{A}_{mn} = c^\dagger_n c_m + c^\dagger_m c_n,
\qquad
\hat{B}_{mn} = i(c^\dagger_n c_m - c^\dagger_m c_n).
\end{equation}
These are the operator versions of $2C^R_{mn}$ and $-2C^I_{mn}$, respectively: $\langle \hat{A}_{mn} \rangle = 2 C^R_{mn}$ and $\langle \hat{B}_{mn} \rangle = -2 C^I_{mn}$.

From the fermionic canonical anticommutation relations $\{c_m, c^\dagger_n\} = \delta_{mn}$, the commutator of $\hat{A}$ and $\hat{B}$ evaluates exactly to
\begin{equation}
[\hat{A}_{mn}, \hat{B}_{mn}] = -2i(\hat{n}_n - \hat{n}_m).
\end{equation}

Substituting into the Robertson uncertainty relation $\Delta A \cdot \Delta B \ge \frac{1}{2} |\langle [A,B] \rangle|$ gives
\begin{equation}
\Delta \hat{A}_{mn} \cdot \Delta \hat{B}_{mn} \ge |\langle \hat{n}_n \rangle - \langle \hat{n}_m \rangle|.
\end{equation}

Translating to the correlation matrix variables $\Delta C^R_{mn} = \frac{1}{2} \Delta \hat{A}_{mn}$, $\Delta C^I_{mn} = \frac{1}{2} \Delta \hat{B}_{mn}$ yields
\begin{equation}
\boxed{\Delta C^R_{mn} \cdot \Delta C^I_{mn} \ge \frac{1}{4} \bigl| \langle n_m \rangle - \langle n_n \rangle \bigr|}. \label{eq:predA}
\end{equation}
This is Prediction~A. The derivation is mathematically exact; no approximation has been introduced beyond the standard Robertson inequality.

\subsection{Physical meaning and novelty}

The Robertson uncertainty relation $\Delta A \cdot \Delta B \ge \frac{1}{2} |\langle [A,B] \rangle|$ is a mathematical theorem valid for any pair of Hermitian operators. For generic operator pairs, however, the commutator expectation $\langle [A,B] \rangle$ is a formal expression that requires knowledge of the full density matrix $\rho$ and provides no operational shortcut. The bound remains a theory-internal consistency condition rather than a relation among independently measurable quantities.

What distinguishes the operator pair $(\hat{A}_{mn}, \hat{B}_{mn})$ is that their commutator evaluates to a simple, directly measurable physical quantity: the difference in occupation numbers between two lattice sites. The bound $\frac{1}{4} |\langle n_m \rangle - \langle n_n \rangle|$ is therefore not merely a formal expression---it is a relation among three independently measurable quantities $(\Delta C^R, \Delta C^I, |\Delta n|)$. This transforms the Robertson inequality from a theory-internal identity into an experimentally testable physical relation.

The content of Prediction~A is the following. In a uniform system where $|\Delta n| \approx 0$, the bound is consistent with zero: classical behavior (arbitrarily small uncertainty product) is permitted by the inequality. In a region with a finite density gradient $|\Delta n| \gg 0$, the bound is strictly positive, enforcing a minimum quantum uncertainty. This ``uncertainty switch''---from classically tolerant to quantum-mandated as the density gradient increases---is the experimentally testable content of Prediction~A. It is the first uncertainty relation whose lower bound is set by a directly observable macroscopic quantity (a density gradient) rather than by a formal commutator expectation that must be computed from the full quantum state.

\subsection{Experimental test in cold-atom quantum gas microscopes}

The prediction can be tested using existing cold-atom quantum gas microscope data. The experiment of Mazurenko et al.~\cite{mazurenko2017} realized the two-dimensional Fermi-Hubbard model with $^6$Li atoms at Hubbard parameters $U/t = 7.4(6)$, $T/t = 0.25(2)$, and doping $\delta = 0.15$ away from half-filling. The harmonic trapping potential creates a natural density gradient from the trap center to the edge.

We select a pair of adjacent lattice sites $(m,n)$ straddling this gradient: site $m$ (closer to trap center) with $\langle n_m \rangle = 0.82(3)$, and site $n$ (further from center) with $\langle n_n \rangle = 0.58(3)$. The density difference is $|\Delta n| = 0.24(4)$, giving a Prediction~A lower bound of $\frac{1}{4} |\Delta n| = 0.060(10)$. For comparison, in the trap center (near-uniform region) the densities are $\langle n_m \rangle \approx 0.85(3)$ and $\langle n_n \rangle \approx 0.84(3)$, yielding $|\Delta n| \approx 0.01(4)$ and a bound consistent with zero: $0.003(10)$. The contrast between the two regions provides a clean experimental signature.

We note that the density values quoted are representative estimates based on the density profiles reported in Mazurenko et al.~(2017), Fig.~2. Precise site-resolved values for the specific adjacent lattice sites in the gradient region would need to be extracted from the raw experimental data in collaboration with the original authors. The uncertainty estimates ($\pm 0.03$) reflect typical single-site imaging fidelity ($\sim 95\%$) at approximately 500 shots.

The quantities $\Delta C^R$ and $\Delta C^I$ can be extracted from shot-to-shot fluctuations in the correlation matrix elements measured via the phase microscope technique~\cite{bruggenjurgen2024} or the kinetic-energy readout protocol of Karch et al.~\cite{karch2024}. Estimated measurement precision: density measurement $\sim 3\%$ per site (500 shots, $95\%$ imaging fidelity); off-diagonal correlation $\sim 5\%$ precision (1000 shots, after phase-to-density mapping); combined uncertainty product precision $\sim 8\%$; signal-to-noise ratio SNR $\sim 12$. This SNR estimate assumes uncorrelated errors between the $\Delta C^R$ and $\Delta C^I$ measurements, which is optimistic. Correlated systematic errors arising from the phase-to-density mapping calibration would increase the effective uncertainty. A dedicated experimental error budget analysis, beyond the scope of this theoretical paper, is needed for a definitive test.

\textbf{Falsification condition.} If a measurement in the gradient region finds $\Delta C^R \cdot \Delta C^I < 0.060$ at $> 95\%$ confidence, Prediction~A is falsified. Such a result would indicate a violation of the Robertson inequality for the pair $(\hat{A}_{mn}, \hat{B}_{mn})$, which would imply either a breakdown of the canonical anticommutation relations or a failure of the standard uncertainty principle in this context---either outcome would be of profound significance.

%% =================================================================
\section{Prediction B: Non-Diagonal Frequency Correlations in Analogue Black Holes}
%% =================================================================

\subsection{Motivation}

Theorem~1 establishes that $C^I$ is a dynamical degree of freedom that cannot be eliminated within a real-valued description. Identity~1 shows that $C^I$ is coupled to $C^R$ through the Hamiltonian $H$, and this coupling generates a perpetual exchange between the two components. In standard quantum dynamics, this exchange is mediated by the system Hamiltonian, and the total symplectic volume $\Omega = \sum dC^R \wedge dC^I$ is conserved.

At a causal horizon---whether an event horizon in general relativity or an analogue sonic horizon in a Bose-Einstein condensate---a qualitatively new situation arises. The $C^I$ that flows toward the horizon cannot be matched by a compensating $C^R$ flow emerging from behind the horizon, because causal disconnection prevents the return coupling. The canonical conjugate cycle is disrupted. The information carried by $C^I$ must therefore manifest itself elsewhere in the accessible region. In an analogue black hole, we propose that this manifests as non-diagonal correlations between different frequency modes of the Hawking radiation.

\subsection{Derivation sketch}

Consider a BEC analogue black hole. The order parameter $\Psi(x,t) = \sqrt{n(x,t)} e^{i\phi(x,t)}$ satisfies the Gross-Pitaevskii equation. Bogoliubov excitations around the condensate are described by mode functions $u_\omega(x)$, $v_\omega(x)$ satisfying the Bogoliubov-de Gennes (BdG) equations
\begin{equation}
\begin{pmatrix}
\mathcal{L} & \mathcal{M} \\
-\mathcal{M}^* & -\mathcal{L}^*
\end{pmatrix}
\begin{pmatrix} u_\omega \\ v_\omega \end{pmatrix}
= \omega \begin{pmatrix} u_\omega \\ v_\omega \end{pmatrix},
\end{equation}
where $\mathcal{L} = -\frac{\hbar^2}{2m}\nabla^2 + V(x) + 2g n(x) - \mu$ and $\mathcal{M} = g n(x)$.

At a sonic horizon where the flow speed exceeds the local sound speed $c_s(x)$, the BdG mode solutions separate into three families: left-moving modes trapped inside the horizon, right-moving modes outside, and modes that convert between positive- and negative-norm branches at the horizon. The scattering matrix $S(\omega, \omega')$ connecting these mode families acquires off-diagonal elements when the sound speed gradient $\partial_x c_s(x)$ is finite.

The density-density correlation function $\langle \delta\hat{n}(x) \delta\hat{n}(x') \rangle$, expressed in the BdG basis and Fourier-transformed to frequency space, contains the off-diagonal contribution
\begin{equation}
C_J(\omega, \omega') \equiv \langle \delta\hat{n}(\omega) \delta\hat{n}(\omega') \rangle_{\omega \neq \omega'}
\propto \int dx\, S^*(\omega, x) S(\omega', x)\, |\partial_x c_s(x)|,
\end{equation}
where $\delta\hat{n}$ is the density fluctuation operator and $c_s(x)$ is the local sound speed. Modeling the sound speed variation as $\delta c_s(x)/c_s \approx \delta\xi/\ell$ across the horizon healing length $\delta\xi$, and approximating the scattering matrix near the horizon by a single-pole form (the horizon acts as a Lorentzian filter with width $\kappa$, the surface gravity), we obtain
\begin{equation}
\boxed{C_J(\omega, \omega') \approx \sqrt{\bar{n}(\omega) \bar{n}(\omega')}\;
\frac{\delta\xi}{\ell}\; \frac{\kappa^2}{(\omega - \omega')^2 + \kappa^2}}, \label{eq:predB}
\end{equation}
where $\bar{n}(\omega) = (e^{\hbar\omega/k_B T} - 1)^{-1}$ is the mean occupation number at the Hawking temperature $T = \hbar\kappa/2\pi k_B$, $\ell$ is the system size, and $\delta\xi$ is the scale over which $c_s(x)$ varies.

\textbf{Approximation status.} This derivation uses two significant approximations: (i) the background occupation $\bar{n}(\omega)$ is taken as uniform (ignoring spatial variation of the mode populations), and (ii) the inhomogeneity is treated perturbatively ($\delta\xi/\ell \ll 1$). The Lorentzian functional form is robust: it follows from the near-horizon scattering pole structure of the BdG equations, which is independent of the detailed density profile. The prefactor $\delta\xi/\ell$ and the signal magnitude estimate of $2$--$5\%$ of $\bar{n}$ near $\omega \approx \omega'$ are order-of-magnitude estimates. A complete quantitative prediction requires full numerical solution of the BdG equations with a spatially varying sound speed profile.

\subsection{Conjecture status and falsification}

We emphasize that Prediction~B is qualitatively different from Prediction~A. Prediction~A follows rigorously from the Robertson uncertainty relation and the fermionic anticommutation relations; its mathematical form is exact. Prediction~B, by contrast, relies on the physical assumption that the disruption of the canonical conjugate cycle at the horizon generates non-diagonal Bogoliubov correlations in the Hawking radiation. This assumption, while physically motivated by the $\hat{J}_{\text{FCS}}$ framework (see Sec.~\ref{sec:discussion}), has not been verified by complete BdG numerics. We therefore classify Prediction~B as a \textbf{conjecture}.

\textbf{Falsification condition.} In a BEC analogue black hole experiment of the Steinhauer type~\cite{steinhauer2016}, measure the density-density correlation function $\langle \delta n(\omega) \delta n(\omega') \rangle$ for $\omega \neq \omega'$ at frequencies near the Hawking temperature. If this quantity is zero within experimental precision (SNR $< 1$), Prediction~B is falsified. The predicted signal magnitude of $2$--$5\%$ is within the sensitivity range of current experiments (estimated SNR $2$--$3$ for a $10$--minute integration time). Falsifying Prediction~B would not affect Theorem~1 or Identity~1, which are mathematical results independent of any physical system. A dedicated numerical study of the BdG equations is in preparation~\cite{huang2026}.

%% =================================================================
\section{Prediction C: Closed-Form Single-Spin-Device Engine Power}
%% =================================================================

\subsection{Model and exact formula}

Consider a three-level quantum system coupled to two thermal baths: a cold bath at temperature $T_C$ coupled to the $|1\rangle \leftrightarrow |2\rangle$ transition, and a hot bath at temperature $T_H$ coupled to the $|1\rangle \leftrightarrow |3\rangle$ transition. The system Hamiltonian is $H = \sum_{i=1}^3 E_i |i\rangle\langle i|$ with energy differences $\hbar\omega_{12} = E_2 - E_1$, $\hbar\omega_{13} = E_3 - E_1$, and $\hbar\omega_{23} = E_3 - E_2 = \hbar(\omega_{13} - \omega_{12})$.

Introducing the dimensionless variables
\begin{equation}
u = \frac{\hbar\omega_{12}}{k_B T_C}, \qquad
v = \frac{\hbar\omega_{13}}{k_B T_H},
\end{equation}
and the Bose occupation factors
\begin{equation}
n_H = \frac{1}{e^{v} - 1}, \qquad
n_C = \frac{1}{e^{u} - 1},
\end{equation}
the relaxation rate of the system to its nonequilibrium steady state is
\begin{equation}
\tau_{\text{relax}}^{-1} = \gamma_H (2n_H + 1) + \gamma_C (2n_C + 1),
\end{equation}
where $\gamma_H$ and $\gamma_C$ are the coupling strengths to the hot and cold baths, respectively.

The ergotropy stored in the steady state $\rho_{ss}$ is $\mathcal{E}_{ss} = \operatorname{Tr}[\rho_{ss} H] - \operatorname{Tr}[H_{\rho_{ss}}^{\text{pass}}]$. Evaluating this for the three-level system with the two-bath coupling yields the closed form
\begin{equation}
\boxed{\mathcal{E}_{ss} = \hbar\omega_{23}\; \frac{e^{u-v} - 1}{e^{u} + 1 + e^{u-v}}}, \label{eq:Ess}
\end{equation}
and the steady-state power is $P = \mathcal{E}_{ss} / \tau_{\text{relax}}$.

The condition for positive power output (engine operation rather than refrigeration) is
\begin{equation}
\frac{T_H}{T_C} > \frac{\omega_{13}}{\omega_{12}}. \label{eq:threshold}
\end{equation}
This threshold has a simple physical interpretation: the hot bath must provide enough energy per excitation to overcome the cold-bath-induced dissipation on the $|1\rangle \leftrightarrow |2\rangle$ transition.

\subsection{Numerical verification}

To verify Eq.~(\ref{eq:Ess}), we performed a numerical scan over 405 parameter combinations spanning the full operating range: $T_H/T_C \in [1.5, 20]$ (15 values), $\omega_{23}/\omega_{13} \in [0.1, 0.9]$ (9 values), $\gamma_H/\gamma_C \in [0.1, 10]$ (3 values). For each parameter point, we diagonalized the Lindblad master equation to obtain the nonequilibrium steady state $\rho_{ss}$ numerically. The analytical formula Eq.~(\ref{eq:Ess}) agrees with the numerical results to within $2\%$ across all 405 points. The small residual deviation arises from the numerical Lindblad steady-state solver's convergence tolerance, not from any approximation in Eq.~(\ref{eq:Ess}). The Carnot bound $\eta \le 1 - T_C/T_H$ is satisfied at every tested parameter point, confirming thermodynamic consistency.

\subsection{Connection to the main framework}

Prediction~C relates to the central results of this paper through Theorem~2. The ergotropy rate $\dot{\mathcal{E}}$ of the three-level engine is driven by $C^I_{mn}$, the imaginary part of the correlation matrix between levels $m$ and $n$. The steady-state ergotropy $\mathcal{E}_{ss}$ given by Eq.~(\ref{eq:Ess}) is the integral of $\dot{\mathcal{E}}$ over the system's relaxation to equilibrium---the total extractable work accumulated by the $C^I$-driven flow before the steady state is reached.

\begin{figure}[t]
\includegraphics[width=\columnwidth]{../latex/Fig1_SSD_phase_diagram.png}
\caption{Steady-state ergotropy $\mathcal{E}_{ss}/\hbar\omega_{13}$ of the SSD quantum engine as a function of the temperature ratio $T_H/T_C$ and frequency ratio $\omega_{23}/\omega_{13}$. The solid line marks the analytical engine--refrigerator boundary, Eq.~(\ref{eq:threshold}). Positive ergotropy occupies 73\% of the scanned parameter space. The analytical expression Eq.~(\ref{eq:Ess}) agrees with numerical Lindblad steady-state diagonalization to within 2\% at every tested point.}
\label{fig:phase}
\end{figure}

\begin{figure}[t]
\includegraphics[width=\columnwidth]{../latex/Fig2_BM_vs_RC.png}
\caption{Strong-coupling analysis. (a) Ergotropy versus coupling strength $\alpha$: Born-Markov baseline (blue), reaction-coordinate numerical (red), and leading-order level-shift estimate (green). (b) Enhancement factor $\mathcal{F}(\alpha)$: the level-shift estimate predicts $\mathcal{F} > 1$ by accounting for the lowered inversion threshold; the full RC numerics yield $\mathcal{F} \approx 0.1$--$0.2$, confirming that population dilution and additional dissipation overwhelm the leading-order enhancement.}
\label{fig:rc}
\end{figure}

%% =================================================================
\section{Discussion}
\label{sec:discussion}
%% =================================================================

We have shown that the fermionic correlation matrix $C_{mn} = \langle c^\dagger_n c_m \rangle$ decomposes into a canonical conjugate pair $(C^R, C^I)$ evolving under coupled dynamics (Identity~1), that this dynamics forces a transition from $\mathbb{R}$ to $\mathbb{C}$ as the minimal sufficient description (Theorem~1), and that the $C^I$ component---precisely the information discarded by the Born rule $|\cdot|^2$ projection---is the unique driver of the ergotropy rate (Theorem~2). We now discuss the broader implications.

\subsection{The Born-Schr\"odinger gap as a resource}

The central message of this work is that $\dot{\mathcal{E}} \propto C^I$: the rate of change of extractable work is driven by the imaginary part of the correlation matrix. A purely mixed state ($C^I = 0$) is a dead battery, regardless of its energy content. A coherent state ($C^I \neq 0$) can deliver power even at the same energy expectation value. The Born rule projection $|\cdot|^2 : \mathbb{C} \to \mathbb{R}_{\ge 0}$ discards $C^I$ in converting the complex correlation matrix into real-valued occupation probabilities. Theorem~2 shows that what is discarded is precisely the quantity that drives work extraction.

This reframes the historical ``gap'' between the Schr\"odinger equation and the Born rule. Far from being a conceptual defect of quantum mechanics, the gap is a necessary structural feature that enables quantum thermodynamic advantage. The practical implication for quantum engineering is significant: strategies that focus solely on decoherence suppression (keeping $C^R$ stable) may miss the more important objective of maintaining $C^I$ (coherent coupling between levels). Device optimization should aim to maximize $C^I$---the imaginary-world current---rather than merely to minimize decoherence.

\subsection{Relation to other work}

\textbf{Goyal, Knuth, and Skilling (2010).} Goyal et al.~\cite{goyal2010} derived complex probability amplitudes from operational symmetry principles: starting from Feynman's sum and product rules for probabilities, they showed that continuous reversible transformations between measurement outcomes force the probability amplitudes to take values in $\mathbb{C}$. Their derivation establishes the necessity of complex numbers at the level of \textit{measurement postulates}. Our work identifies a complementary mechanism at the level of \textit{correlation matrix dynamics}: the canonical conjugate structure of $C^R$ and $C^I$ forces an escape from $\mathbb{R}$ to $\mathbb{C}$ for purely algebraic reasons related to diagonalizability. The two derivations reach the same conclusion---complex numbers are necessary, not merely convenient---through independent routes.

\textbf{Renou et al. (2021).} Renou et al.~\cite{renou2021} demonstrated that real quantum mechanics makes statistically distinguishable predictions from complex quantum mechanics in network Bell scenarios. Their result operates at the \textit{operational level}: it establishes \textit{that} real quantum mechanics fails as an empirical theory in certain scenarios. Our result operates at the \textit{structural level}: it shows \textit{why} real quantum mechanics must fail, by tracing the failure to the algebraic incompleteness of $\mathbb{R}$ under the canonical conjugate dynamics of the correlation matrix. The two results are complementary.

\textbf{Orion et al. (2026).} Orion et al.~\cite{orion2026} established topological sum rules connecting geometric phases to Hamiltonian topology. Their work addresses the classification of quantum states by topological invariants. Our work addresses the complementary question of dynamics: given a topological class, what drives the flow of extractable work? The two frameworks connect through the symplectic structure $\Omega = \sum dC^R \wedge dC^I$, which is related to the Berry curvature underlying the topological sum rules.

\textbf{Structural necessity across quantum foundations.} The present argument---that $\mathbb{C}$ is forced by the spectral structure of quantum mechanics, not merely convenient for it---is corroborated by independent lines of evidence across quantum foundations. Wigner's theorem identifies time reversal with complex conjugation~\cite{wigner1931}, making the antiunitary sector of quantum symmetries intrinsically complex. The spin-statistics theorem relies on analytic continuation to the complex Lorentz group~\cite{streater1964}. Moretti and Oppio~\cite{moretti2017} proved that Poincar\'e symmetry forces a unique complex structure on the real quantum-mechanical state space. Niestegge~\cite{niestegge2020} showed that local tomography rules out both real and quaternionic quantum mechanics, leaving $\mathbb{C}$ as the only consistent option. These results, together with the operational falsification of real quantum mechanics by Renou et al.~\cite{renou2021} and the algebraic reconstruction of complex amplitudes by Goyal et al.~\cite{goyal2010}, form a converging body of evidence: complex numbers are not an optional convention but a structural requirement of quantum theory. Our result adds the dynamical mechanism---canonical conjugate dynamics of the correlation matrix cannot be spectrally decomposed over $\mathbb{R}$---to this converging picture.

\subsection{The Maslov index and the discrete sector of $C^I$}

The Maslov index $\mu$ is a topological invariant that counts the number of sign changes of the symplectic $2$-form $\Omega$ along a closed trajectory in phase space. In the context of the canonical conjugate pair $(C^R, C^I)$, the Maslov index acquires a direct physical meaning: it counts the number of times the trajectory in the $(C^R, C^I)$ plane crosses the $C^I = 0$ axis, which corresponds to a phase jump of $\Delta\phi = -(\pi/2)\mu$ in the complex correlation $C = C^R + iC^I$. The integer-valuedness of $\mu$ for closed orbits implies that certain phase shifts are topologically protected: they arise from the geometry of the symplectic flow rather than from the specific Hamiltonian. This observation connects to the well-known distinction in fault-tolerant quantum computing between Clifford gates (which correspond to integer $\mu$) and the $T$ gate (which requires $\mu = 1/2$ and therefore demands active error correction).

Within the framework of the present paper, the Maslov index is the discrete, topological sector of $C^I$---the part that is protected by the symplectic geometry of the correlation matrix dynamics rather than by continuous dynamical symmetries. The conservation of $\Omega$ under the Heisenberg flow (Identity~1) guarantees that this topological content is preserved under unitary evolution.

\subsection{Open problem: an FCS-based information current}

An earlier version of this work introduced a constraint $\hat{J} = dS_{vN}/d\tau + \nabla_\mu J^\mu_G$ aimed at tracking the information discarded by the Born rule. However, the von Neumann entropy $S_{vN} = -\operatorname{Tr}[\rho \ln \rho]$ is itself defined through the spectral decomposition of $\rho$, which presupposes the Born rule projection $|\cdot|^2$. Using $S_{vN}$ to quantify the ``gap'' created by that same projection is circular.

A more consistent approach replaces the von Neumann entropy with the cumulant generating function of Full Counting Statistics (FCS)~\cite{levitov1993,levitov1996}. The FCS cumulant generating function is
\begin{equation}
\chi(\lambda) = \ln \operatorname{Tr}[\rho e^{i\lambda \hat{N}}],
\qquad
C_k = (-i)^k \left. \frac{\partial^k \chi(\lambda)}{\partial \lambda^k} \right|_{\lambda = 0}.
\end{equation}
The second cumulant $C_2 = \langle \hat{N}^2 \rangle - \langle \hat{N} \rangle^2$ can be expressed in terms of the correlation matrix as
\begin{equation}
C_2 = \sum_m \langle n_m \rangle (1 - \langle n_m \rangle) - \sum_{m \neq n} |C_{mn}|^2
= \sum_m \langle n_m \rangle (1 - \langle n_m \rangle) - \sum_{m \neq n} \bigl[ (C^R_{mn})^2 + (C^I_{mn})^2 \bigr].
\end{equation}
The second term directly involves $C^I$ without requiring the spectral decomposition of $\rho$. Crucially, the FCS cumulants are operationally measurable: $C_k$ are the cumulants of the full counting distribution $P(N)$, which can be obtained by simply counting particles in repeated experiments. No diagonalization of $\rho$ is needed.

Define the FCS information current
\begin{equation}
\hat{J}_{\text{FCS}} = \frac{dC_2}{d\tau} + \nabla_\mu \mathcal{J}^\mu_{\text{info}},
\end{equation}
where $\mathcal{J}^\mu_{\text{info}}$ is the FCS information current defined through the continuity equation for the cumulant generating function. In standard quantum systems, $\hat{J}_{\text{FCS}} = 0$: the particle number variance lost by the system equals the variance gained by the environment. At a causal disconnection, the FCS cumulants show anomalous growth or decay---$\hat{J}_{\text{FCS}} \neq 0$ signals a genuine information gap. The complete development of this framework, including the Keldysh formalism treatment and connection to the geometric entropy current of Jacobson~\cite{jacobson1995}, is deferred to future work.

%% =================================================================
%% Acknowledgments
%% =================================================================
\begin{acknowledgments}
[To be completed.]
\end{acknowledgments}

%% =================================================================
%% Bibliography
%% =================================================================
\bibliography{main_pra}

\end{document}
